\section{状態数と写像：数える対象を先に固定する}
\label{theory:discrete:bitcount}
二値の状態を持つn個の位置について、各位置の選択を独立に行えるなら全状態は\(2^n\)通りである。位置数の比と状態数の比は異なる。位置がmからnへ増える場合の倍率は\(2^{n-m}\)であり、\(n/m\)ではない。IPv4の32bitとIPv6の128bitでは、この全ビット列数の比が\(2^{96}\)となる\parencite{rfc4291_addresses}。これは表現可能な全列を数える計算で、実際に利用者へ割り当てるアドレスの割合を推定するものではない。

2進数から10進表記へ変える場合も、位取りと並びの順は別である。IPv4は上位から8bitずつ四塊へ分け、各塊を10進数へ変える。16進や10進の表示は同じビット列の書き方であり、それだけで状態数は変わらない。

\subsection{直積からの写像と積の法則}
\label{theory:discrete:mappings}
有限集合Xの各要素へ、集合Yのいずれかの要素を像として選ぶとする。\(|X|=m\)、\(|Y|=k\)なら、全写像数は\(k^m\)である。各入力に一つの出力を割り当てることが写像の条件であり、異なる入力の像が同じでもよい。単射や全射の条件が付けば数え方は変わる。

\(|A|=a\)、\(|B|=b\)なら直積A\(\times\)Bは\(ab\)個の順序対を持つ。Cの要素数がcなら、直積からCへの写像は\(c^{ab}\)通りとなる。\(abc\)は入力と出力の組の総数であり、全入力への割当てを一つとして数える写像数ではない。

\section{包除原理：重複の次数で符号を交替する}
\label{theory:discrete:inclusion}
既存の第\ref{theory:information:sets}節の二集合の式を多集合へ拡張する。各集合の人数を単純に足すと、複数の集合へ属する人を重複して数える。二重の交差を引き、三重を足し、四重を引くという補正を続ければ、和集合の要素を各一回にできる\parencite{mit_discrete2015}。
\[
\left|\bigcup_{i=1}^{m}A_i\right|
=\sum_{\varnothing\ne I\subseteq\{1,\ldots,m\}}
(-1)^{|I|+1}\left|\bigcap_{i\in I}A_i\right|.
\]
ある要素がk個の集合に属すれば、その係数は\(\sum_{j=1}^k(-1)^{j+1}\binom{k}{j}=1\)となる。これは\((1-1)^k=0\)から得られ、重複補正の根拠となる。

四集合の単集合・任意の二重・任意の三重交差の数がそれぞれ一定なら、二重には\(\binom42=6\)、三重には\(\binom43=4\)を掛ける。交差の数が等しいことから交差が一つだけあるとは考えない。補集合を問う場合は、包除で得た和集合数を全体集合数から引く。原典で上線がどこまで付くかを確認する。

\section{算術木と後置記法：値より先に構造を読む}
\label{theory:discrete:expression}
算術式は葉に演算数、内部節点に演算子を置いた順序付き二分木で表せる。減算や除算は交換できないため、左右の子を区別する。木が与えられたら、左部分式を求め、右部分式を求め、根の演算を作用させる。中置表記では必要な括弧でこの構造を保持する。

後置記法は、左部分木・右部分木・根の順に記号を出す後行順走査である。例えば\(a(b-c)\)は\texttt{abc-*}となる。二項演算では、演算子が現れた時点でスタックから先に右の値、次に左の値を取り出す。\texttt{ab-}を\(b-a\)としてはいけない。

\begin{tcolorbox}[title=情報工学への接続]
式の評価、構文解析、再帰呼出しの管理では、対象の木構造と走査の順序を対応させる。この接続は学習上の補足であり、算術木の一問から全ての構文解析法の出題頻度を推定しない。
\end{tcolorbox}

\subsection{スタックの状態と最後の戻り値}
\label{theory:discrete:stack}
既存のスタック節で定義したPUSHとPOPは、要素の値と順序を保って追跡する。底から\([1,2,3]\)があるとき、POPは3を返して\([1,2]\)を残す。最後に返った値と最後に残った値は同じとは限らない。途中で別の値をPUSHしたら、その値が先に取り出される。空へのPOP等は別の仕様を要するが、ここでの過去問の操作列には現れない。

\section{構文経路：通る順序と反復する単位を読む}
\label{theory:discrete:syntax}
構文図は開始から終了までの有向経路に沿って記号を消費する。英字一つの後に数字を1個以上読む単位を0回以上繰り返し、最後に英字一つを読む図なら、受理する文字列は\((LD^+)^*L\)である。ここでLは英字、Dは数字を表す。開始側の上の迂回は反復0回を許すので、最後の英字一つだけも受理する。

反復の戻り先を確認する。この図では数字の直後から数字の手前へ戻る小さなループがあり、連続数字を許す。より外側のループは英字と1個以上の数字をまとめて繰り返す。数字を経ずに英字同士を連続させる経路はない。図に存在しない経路を、通常の識別子の慣習から足してはいけない。第\ref{theory:information:grammar}節の再帰規則と同様に、各候補の全記号を読み終えて終了点へ届くかを検査する。

\section{二分探索木：挿入順を祖先の先行条件へ変える}
\label{theory:discrete:bst}
二分探索木は、ある節点より左の部分木へ小さいキー、右へ大きいキーを置く。挿入では根から比較を繰り返し、空のリンクへ新節点を置く。木の形は挿入順に依存する\parencite{princeton_bst}。

同じ木を作る順番を考えるには、各部分木の根を、その子孫より先に挿入する条件へ変える。異なる枝の挿入は交互に行ってよいが、親より先にその子となるはずのキーを置くと、そのキーが部分木の根になってしまう。平均的な探索効率の議論は入力順の仮定を必要とし、全挿入順で対数時間になるという保証ではない。

\section{走査の不変条件と分岐：成功例と失敗例の双方を見る}
\label{theory:discrete:invariants}
アルゴリズムの正しさは、途中状態が何を表すかを定めて考える。最小値走査では、i個を読んだ時点でminがそのi個の最小であることを不変条件にできる。初期値が全入力より大きく、更新で\(\min(\mathrm{min},a_i)\)を保てば、最後の正しさを示せる。

最大値の更新を最小値更新のelseへ置くと、最小を更新した入力では最大の検査をしない。最初の値が最大となる列では最大を取りこぼすが、後の入力がより大きい列では正しい最大を得ることもある。「失敗する場合がある」と「常に失敗する」を、反例と成功例を分けて判断する。

試し割りによる素数判定では、\(N=uv\)なら片方の因子が\(\sqrt N\)以下である。2から平方根まで全てを調べて割り切れなければ素数である。\(I\le\sqrt N\)で継続し、\(N\bmod I=0\)なら非素数へ進む。平方数の境界を含むこと、割る数と割られる数、真偽の矢印を確認する。

\section{拡張互除法：同じ最大公約数でも係数は一意でない}
\label{theory:discrete:extended}
\(a=qb+r\)なら\(\gcd(a,b)=\gcd(b,r)\)である。余りの計算と同時に、各余りを最初のa,bの整数線形結合として保存すれば、最後の非零余りdを\(ax+by=d\)と表せる\parencite{mit_discrete2015}。
\[
\begin{pmatrix}b\\r\end{pmatrix}
=T(q)\begin{pmatrix}a\\b\end{pmatrix},
\qquad
T(q)=\begin{pmatrix}0&1\\1&-q\end{pmatrix}.
\]
次の操作は左から行列を掛ける。商が\(q_1,q_2,q_3\)なら合成は\(T(q_3)T(q_2)T(q_1)\)であり、行列を逆の順に並べない。指定行列の第二行が最後の非零余りを与えるとき、その行を係数として読む。

\(ax+by=d\)を満たす一組があれば、\(x'=x+(b/d)t\)、\(y'=y-(a/d)t\)も任意の整数tについて同じ等式を満たす。したがって候補を線形和へ代入するだけでは複数が合う場合がある。過去問で「指定行列の行」を求める条件があるなら、その行列積も照合して一つを選ぶ。

\section{漸近上界：定数・有限区間・増大率を分ける}
\label{theory:discrete:asymptotic}
\(f(n)=O(g(n))\)は、ある正の定数cと境界\(n_0\)以後で\(f(n)\le cg(n)\)となることを表す。非負関数の上界であり、常に同じ値になることや最小次数を指定する等号ではない\parencite{mit_discrete2015}。

比を取ると、\(\log n/n^\varepsilon\to0\)、\(n^k/a^n\to0\)が\(\varepsilon>0,a>1\)の条件で成り立つ。指数同士なら底の比が重要で、\(n^3(3/4)^n\to0\)である。二重指数と大きなべきの指数を比較するときは対数を取る。\(2^{2^n}/10^{n^{100}}\)の自然対数は\(2^n\ln2-n^{100}\ln10\)で、十分大きいnでは正無限大へ発散する。有限の小さなnで後者が大きく見えることは、漸近上界の成立を保証しない。

\section{入れ子の再帰：基底から低い階層を先に得る}
\label{theory:discrete:recursion}
再帰関数は、最小の場合の値と、その値へ戻る規則を対応させて読む。アッカーマン関数の\(A(0,n)=n+1\)から、\(A(1,0)=2\)、\(A(1,n+1)=A(1,n)+1\)となり\(A(1,n)=n+2\)を得る。これを用いれば、\(A(2,0)=3\)、\(A(2,n+1)=A(2,n)+2\)、従って\(A(2,n)=2n+3\)となる。

入れ子の\(A(m,A(m+1,n))\)では、内側の値を外側の引数へ代入する。引数の添字が一つ減ることだけを数えて値を足さない。低い階層の式を導いてから必要な値を求めると、再帰を同じ形で何度も書く負荷を減らせる。
